\chapter{Research Case Studies and Take-Homes}\label{ch:iv:research-case-studies-and-take-homes}

The e-mail arrives at nine: a file of 156\,040 minute records of quotes and trades for twenty stocks, one sentence (``Is
there anything predictable here?'') and a four-hour deadline. The submission that wins is not the one with the most models.
It is the one whose first page says what was checked, what was found, how sure the author is, and what she would do with
another week. A \omterm{def:iv:research-case-studies-and-take-homes:takehome}{research case study} is the closest an interview comes to the job itself, and it is scored the way a research
review scores a result (One Quant Book~7, chapter~1): on the question asked, the honesty of the test and the clarity of the
conclusion.

\section{What a reviewer reads first: the memo}

\begin{definition}[Take-home assignment, research case study]\label{def:iv:research-case-studies-and-take-homes:takehome}
A \emph{take-home assignment}\index{take-home assignment} is an interview task done by the candidate outside the interview,
under a deadline of hours or days, with a dataset or a specification and a written deliverable. A \emph{research case
study}\index{research case study} is a take-home or in-person exercise in which the candidate investigates an open question on
supplied data and presents a conclusion, and is assessed on the investigation as much as on the result.
\end{definition}

Reviewers read the first page and decide how much of the rest to read. The first page must stand alone.

\begin{method}[The one-page memo]\label{met:iv:research-case-studies-and-take-homes:memo}
\begin{enumerate}
\item \emph{Question}: the one precise question investigated, and why it was chosen.
\item \emph{Data}: what was checked and fixed (duplicates, crossed quotes, corporate actions, gaps, time zones), with counts.
\item \emph{Result}: the estimate with its standard error, and one chart.
\item \emph{Robustness}: the two or three checks that could have killed it, and what they showed.
\item \emph{So what}: whether it is economically meaningful after costs, and what you would do next.
\item \emph{Limits}: what you did not do and why. Code in an appendix that runs from a clean checkout.
\end{enumerate}
\end{method}

\section{The four hours}

Most candidates spend the time on models; most failures come from data. A plan that works: thirty minutes on data checks,
thirty on choosing one clean question, ninety on one honest test with its robustness checks, forty-five on the memo, fifteen
on re-running the code from scratch. A second question is started only if the first is finished.

\begin{example}[The data checks on the hook's file]\label{ex:iv:research-case-studies-and-take-homes:checks}
The chapter's generated dataset (\texttt{iv\_case.py}) contains three planted defects that the standard checks find: 40
duplicated (stock, day, minute) records, 15 rows with the bid above the ask (swapped fields), and a jump of about 50\% in one
stock's price between two days that is a 2-for-1 split not adjusted in the data. Each is fixed and counted in the memo. A
candidate who computes daily returns without finding the split reports a $-50\%$ day.
\end{example}

\section{Two worked case studies with a reviewer's comments}

\emph{Case one: does order-book imbalance predict the next minute?} The question chosen from ``anything predictable'':
the slope of the next minute's mid-price return on the current imbalance, pooled over stocks. On the cleaned data the slope is
0.48 basis points per unit of imbalance with a standard error of 0.022 clustered by stock-day, a $t$-statistic near 22, and the
decile means line up with it (\cref{fig:iv:research-case-studies-and-take-homes:case}). The memo then asks what the effect is
worth: the strongest decile predicts about 0.45 basis points, against a round trip of 6 basis points of spread.

\begin{omfigure}
\begin{tikzpicture}
\begin{axis}[width=10cm,height=5.2cm,xlabel={\small order-book imbalance (decile centre)},
  ylabel={\small next-minute mid return (bp)},xmin=-1,xmax=1,ymin=-0.6,ymax=0.6,xtick={-1,-0.5,0,0.5,1},
  tick label style={font=\footnotesize},grid=major,grid style={gray!20},table/col sep=comma,
  legend cell align=left,legend style={font=\scriptsize,at={(0.5,-0.32)},anchor=north,
  /tikz/every even column/.append style={column sep=8pt}},legend columns=2]
\addplot[only marks,mark=*,mark size=1.6pt,omDef,error bars/.cd,y dir=both,y explicit]
  table[x=centre,y=mean,y error expr=2*\thisrow{se}]{figdata/interviews/20-research-case-studies-and-take-homes/case.csv};
\addplot[omThm,thick,dashed] table[x=centre,y=planted]{figdata/interviews/20-research-case-studies-and-take-homes/case.csv};
\legend{decile mean $\pm$ 2 standard errors,planted relation (0.5 bp per unit)}
\end{axis}
\end{tikzpicture}

\omcaption{The take-home's first case: the next minute's mid-price return by decile of current imbalance, on the cleaned
generated dataset of twenty stocks and twenty days, with the relation planted by the generator. Statistically overwhelming,
economically a fraction of a basis point. Data: \texttt{fig\_iv\_case.py}.}\label{fig:iv:research-case-studies-and-take-homes:case}
\end{omfigure}

\begin{remark}[Reviewer's comments on case one]\label{rem:iv:research-case-studies-and-take-homes:rev1}
Strong: the data section counts and fixes the three defects; the returns are computed from mid prices, not trades, and the
memo explains why (trade prices bounce between bid and ask, which gives trade-price returns a lag-one autocorrelation of about
$-0.21$ here, against zero for mid prices; Roll's model of the bounce, One Quant Book~4, chapter~21); the standard error is
clustered. Strongest of all: the ``so what'' paragraph says the effect cannot pay a crossing strategy and asks whether it
could improve the timing of passive orders instead. Missing: a split of the sample in time, and a check that the effect is
not driven by one stock.
\end{remark}

\emph{Case two: a calendar effect.} A second candidate tests mean returns by weekday and month, 60 cells, finds one cell with
$t = 3.1$ and writes a memo about ``the March-Tuesday effect''. With 60 independent null cells the chance that at least one
reaches $|t| \ge 3.1$ is about 11\%, and a Bonferroni threshold for 60 cells at 5\% is $|t| \ge 3.34$.

\begin{remark}[Reviewer's comments on case two]\label{rem:iv:research-case-studies-and-take-homes:rev2}
The analysis is competent and the conclusion is wrong: the memo reports the best of 60 tests as if it were the only one. A
good memo says ``I tested 60 cells; one exceeded the nominal threshold, as expected by chance about one time in nine; nothing
survives a multiple-testing correction; I found no calendar effect.'' A null result stated with its power is a good result
(One Quant Book~4, chapter~12).
\end{remark}

\section{Presenting and defending the result}

A take-home is often followed by a presentation in which reviewers challenge the result. The questions are predictable: is it
the bid-ask bounce, is it one stock or one day, is it there out of sample, what does it cost to trade, how many things did you
try. The best defence is to have asked them first and to put the answers on the memo's first page. When a challenge is right,
concede it and quantify its effect; a candidate who revises a conclusion under good questioning scores better than one who
defends a weak one.

\section{Worked answers}

\begin{example}[How much data the question needs]\label{ex:iv:research-case-studies-and-take-homes:power}
\emph{``Before you start: if the effect you are looking for is 0.2 basis points of next-minute return per unit of signal,
the signal has unit variance and the residual standard deviation is 5 basis points, how many independent minutes do you
need to see it at $t = 3$?''} The slope's standard error is about $5/\sqrt{n}$, so $t = 0.2\sqrt n/5 = 3$ needs $\sqrt n =
75$: $n \approx 5\,600$ minutes, about fourteen and a half days of one stock's 390-minute sessions if the minutes were
independent. They are not (minutes of one day share news and liquidity), so the file of twenty stocks and one week is
borderline, and the memo should say so before it reports anything. Computing the sample the question needs, in the first ten
minutes, decides what question the four hours can answer.
\end{example}

\section{Question bank}

\begin{interviewq}[$\star$ \iqroles{researcher, mle} \iqfirm{systematic fund}]\label{iq:iv:research-case-studies-and-take-homes:1}
You receive a file of minute quotes and trades for twenty stocks. List the data checks you run in the first thirty minutes,
in order.
\end{interviewq}

\begin{interviewq}[$\star$ \iqroles{researcher} \iqfirm{any}]\label{iq:iv:research-case-studies-and-take-homes:2}
What goes on the first page of a take-home memo, and what goes in the appendix?
\end{interviewq}

\begin{interviewq}[$\star$ \iqroles{researcher, trader} \iqfirm{market maker}]\label{iq:iv:research-case-studies-and-take-homes:3}
The prompt is ``Is there anything predictable here?''. Turn it into one clean question you can answer in four hours, and say
why you chose it.
\end{interviewq}

\begin{interviewq}[$\star$ \iqroles{researcher, mle} \iqfirm{systematic fund}]\label{iq:iv:research-case-studies-and-take-homes:4}
Your four hours produced no significant effect. How do you write it up?
\end{interviewq}

\begin{interviewq}[$\star\star$ \iqroles{researcher, developer} \iqfirm{systematic fund}]\label{iq:iv:research-case-studies-and-take-homes:5}
In the chapter's dataset, what do the three standard checks find, and how do you fix each defect? What goes wrong if you miss
the split?
\end{interviewq}

\begin{interviewq}[$\star\star$ \iqroles{researcher} \iqfirm{market maker}]\label{iq:iv:research-case-studies-and-take-homes:6}
You regress the next minute's mid return on imbalance and get 0.48 basis points per unit with a clustered standard error of
0.022. Why cluster, and by what? Is the effect real?
\end{interviewq}

\begin{interviewq}[$\star\star$ \iqroles{researcher, trader} \iqfirm{market maker}]\label{iq:iv:research-case-studies-and-take-homes:7}
Returns computed from trade prices show a lag-one autocorrelation of $-0.21$; from mid prices, zero. Explain, and give the
formula that predicts $-0.21$ from a half-spread of 3 basis points and a minute volatility of 5.
\end{interviewq}

\begin{interviewq}[$\star\star$ \iqroles{researcher, mle} \iqfirm{any}]\label{iq:iv:research-case-studies-and-take-homes:8}
How do you split four hours between data, analysis and writing, and what do you leave out?
\end{interviewq}

\begin{interviewq}[$\star\star$ \iqroles{researcher} \iqfirm{multi-manager fund}]\label{iq:iv:research-case-studies-and-take-homes:9}
A reviewer reads two memos on the same data: one fits five models and reports the best $R^2$; one tests one hypothesis with
robustness checks and finds a small effect. Which scores higher, and what would you write in the margin of each?
\end{interviewq}

\begin{interviewq}[$\star\star\star$ \iqroles{researcher, trader} \iqfirm{market maker}]\label{iq:iv:research-case-studies-and-take-homes:10}
The imbalance effect is 0.5 basis points per unit of imbalance and the spread is 6 basis points. Can a strategy that crosses
the spread trade it? If not, how could the effect still be worth money?
\end{interviewq}

\begin{interviewq}[$\star\star\star$ \iqroles{researcher} \iqfirm{systematic fund}]\label{iq:iv:research-case-studies-and-take-homes:11}
A candidate reports a ``March-Tuesday effect'' with $t = 3.1$ after testing 60 weekday-month cells. What is the chance of such a
result with no effect at all, and what threshold would Bonferroni use?
\end{interviewq}

\begin{interviewq}[$\star\star\star$ \iqroles{researcher, mle} \iqfirm{multi-manager fund}]\label{iq:iv:research-case-studies-and-take-homes:12}
In the presentation, a reviewer says: ``It is one stock and one week.'' How do you answer with data in five minutes, and what do
you concede if the reviewer is right?
\end{interviewq}

\begin{omsources}
\item R.~Roll, ``A simple implicit measure of the effective bid-ask spread in an efficient market'', \emph{Journal of Finance}
39(4), 1984.
\item One Quant Book~7, chapters 1, 3 and~20 (the research process, point-in-time data, overfitting); One Quant Book~4, chapters
12 and~21 (multiple testing; high-frequency econometrics).
\end{omsources}
